\scanpage{85}
% archive sheet number: 84
% block text | line 1 heading
\deleted{\unsure{Corollaire}} \uline{Cas particulier}\qquad\uline{$H^1(\Gamma,K^*)=0$}.

% block text | lines 2-4
Ici, il suffit de vérifier qu'on peut trouver $x\in K^*$ tel que $\sum f(\unsure{\gamma|\sigma})\gamma x\ne0$, et l'indépendance \uline{linéaire} des automorphismes est suffisante.

% block text | lines 5-9
\uline{Généralisation} Soit $\underline{O}$ un anneau $\frac12$ local noethérien, et $\Gamma$ un groupe fini opérant \uline{fidèlement} sur $\underline{O}$ et \deleted{par} transitif sur les idéaux maximaux. Soit $\Gamma_d$ le \deleted{sous} groupe de stabilité d'un idéal maximal $\underline{m}$ de $O$, donc $\Gamma_d$ opère dans $O/m=k$.

% block text | lines 10-11
Soit $\Gamma_i$ le noyau de la représentation de $\Gamma_d$ \deleted{\unsure{dans le gr.|dans les}}\inserted{par} automorphismes du corps $O/m$.

% block text | lines 12-13
Soit enfin $A$ une $\underline{O}$-algèbre \deleted{\unsure{finie} \inserted{de type fini}} \deleted{Alors} où \inserted{$\Gamma$ opère tel que} \deleted{\unsure{$a\in A$}} \deleted{L'application naturelle} $\gamma(\lambda a)=(\gamma\lambda)(\gamma a)$. Alors les applications
% block diagram | line 14
\diagram{\begin{tikzcd}[ampersand replacement=\&,column sep=small]
H^1(\Gamma,A^*)\arrow[r]\&H^1(\Gamma_d,A_m^*)\arrow[r]\&H^1(\Gamma_i,A_m^*)\deleted{\unsure{/mA}}\arrow[r,"\deleted{\unsure{\longrightarrow}}" description,draw=none]\&\deleted{\unsure{H^1(\Gamma_i,(A/mA)^*)}}
\end{tikzcd}}
% block text | lines 15-19
sont \uline{injectives}. [N.B. \deleted{\unsure{$H^1(\Gamma_i,A/mA)$} \illegible{2}} dans \unsure{les cas les plus importants|le cas le plus important}, \deleted{\unsure{$A/mA$}} $\Gamma_i$ opère trivialement sur $A/mA$, donc $H^1(\Gamma_i,(A/mA)^*)$ est formé des classes d'homomorphismes -- à aut. intérieurs \inserted{de $A/mA$} près -- de $\Gamma_i$ dans $(A/mA)^*$]

% block text | lines 20-22
Soient $f,g\in Z^1(\Gamma,A)$, supposons qu'elles soient cohomologues dans \deleted{$A_m$}\inserted{$\unsure{H^1(\Gamma_d}$}, i.e. qu'il existe $a,\deleted{\unsure{b}}\in A$, $\lambda\in O\setminus m$, \deleted{\unsure{$b\in A$} \illegible{3}} $\unsure{b'|b}\in A$, $\mu\in O\setminus m$, tels que
% block display | line 23
\[
\deleted{(a\lambda^{-1})(b\mu^{-1})=}\quad\lambda^{-1}\mu^{-1}ab=\inserted{\lambda^{-1}\mu^{-1}ba=1}
\]
% block text | line 23 right and line 24 left
dans $A_m$ i.e. il existe $\nu\in O\setminus m$ tel que
% block display | line 24
\[
\deleted{\unsure{\nu\lambda\mu}}\nu ab=\nu\deleted{\unsure{\lambda\mu}}ba=\nu\lambda\mu\ ;
\]
% block text | line 24 right
et \unsure{les équations}
% block display | bottom insertion
\[
\inserted{g(\gamma)=(\unsure{a|a'}\lambda^{-1})^{-1}f(\gamma)\gamma(a\lambda^{-1})}
\]
% block text | line 25 right
dans $A_m$, i.e.
